\documentclass[10pt,a4paper]{article}
\usepackage{amsmath,amssymb,amsthm,graphicx}
\usepackage{xfrac}
\newenvironment{enumeratei}
{\bgroup\def\theenumi{\roman{enumi}}\def\theenumii{\arabic{enumii}}\begin{enumerate}}
{\end{enumerate}\egroup}

\hyphenation{addi-tion-al belong belongs denote exchan-ge exchan-ged exchan-ges}

\newcommand{\psfrac}[2]{\sfrac{(#1)}{#2}}

\newcommand\mto{\mathchoice{\longmapsto}{\mapsto}{\mapsto}{\mapsto}}
\let\oldbigsqcup\bigsqcup
\renewcommand{\bigsqcup}{\mathop{\textstyle\oldbigsqcup}\displaylimits}

\let\epsilon\varepsilon

\DeclareMathOperator{\Ess}{Ess}
\DeclareMathOperator{\Leb}{Leb}

\usepackage{tikz}
\usetikzlibrary{plotmarks}

\theoremstyle{plain}
\newtheorem{theorem}{Theorem}[section]
\newtheorem{lemma}[theorem]{Lemma}
\newtheorem{proposition}[theorem]{Proposition}
\newtheorem{corollary}[theorem]{Corollary}
\newtheorem{claim}{Claim}

\theoremstyle{definition}
\newtheorem{remark}[theorem]{Remark}
\newtheorem{example}[theorem]{Example}

\newcommand{\A}{\mathcal{A}}
\newcommand{\B}{\mathcal{B}}
\newcommand{\I}{\mathcal{I}}

\newcommand{\Z}{\mathbb{Z}}
\newcommand{\N}{\mathbb{N}}
\newcommand{\R}{\mathbb{R}}
\title{A conjugacy-satisfying interval exchange transformation with distinct continuity lengths is symmetric}
\author{Przemyslaw Berk, Frank Trujillo}
\date{doi:10.5802/jep.302; CC BY4.0}
\begin{document}
\maketitle
\noindent\textit{Editorial note.} The counted statement is the author's labelled lemma environment at source0.tex lines 489-494; its proof environment opens at line 496, the line after the blank separator, so the pairing is by adjacency and the retained proof range 496-508 is the whole environment.\par\medskip
\noindent\textit{Editorial note.} Two numbered environments are retained as context: eq: defofintervalorder (source lines 374-376), which defines the order used at source line 507, and eq: standardconjugacy (source lines 463-465), which the statement itself invokes.\par\medskip
\noindent\textit{Editorial note.} The author label of this lemma is written with a trailing comma in the source; the claim key is recorded in normalised form.\par\medskip
\noindent\textit{Editorial note.} This source is new to the material-authoring work and its preamble is built from the author's own macro and theorem block (source lines 8-45) rather than from a prepared adapter.\par\medskip
\noindent\textit{Editorial note.} This source has no label-number dump and the retained counter values are left empty, so the printed numbering of the card is the wrapper's default; the identity of the counted item is its author label.\par\medskip
\noindent\textit{Editorial scope.} One original result of the pinned publisher article, counted once: a sufficient criterion for an interval exchange transformation to be symmetric (author label lem: condition for symmetric iets).\par\medskip
\noindent\textit{Allowed standard prerequisites (editorial).}\begin{itemize}
\item Non-degenerate interval exchange transformation. An IET given by a permutation and a length vector, with all continuity intervals of positive length.. It is the object of the statement.
\item Conjugacy equation. The reflection of the interval conjugates the transformation to its inverse off the discontinuities.. It is the hypothesis of the statement.
\item Order on intervals. One interval is below another when every point of the first is below every point of the second.. It is the order that the argument shows is reversed both by the transformation and by the reflection.
\end{itemize}\medskip
\begin{equation}\label{eq: defofintervalorder}
J_1<J_2\iff x<y \text{ for any }x\in J_1\text{ and any }y\in J_2.
\end{equation}

\begin{equation}\label{eq: standardconjugacy}
\mathcal{I}_I\circ T(x)=T^{-1}\circ\mathcal{I}_I(x) \quad \text{if } x \neq \partial I_{\alpha} \text{ for any } \alpha \in \A.
\end{equation}

\begin{lemma}
\label{lem: condition for symmetric iets,}
Let $T=(\pi,\lambda): I \to I$ be a non-degenerate IET. If $T$ satisfies
\eqref{eq: standardconjugacy}
and $\lambda_{\alpha}\neq \lambda_{\beta}$ for any distinct $\alpha, \beta \in \mathcal{A}$, then $T$ is symmetric.
\end{lemma}

\begin{proof}
Arrange the intervals $\{I_{\alpha}\}_{\alpha\in \A}$ according to their lengths, such that $\lambda_{\alpha_1}>\lambda_{\alpha_2}>\dots> \lambda_{\alpha_d}$.
Note that $\I_I\circ T$ is continuous on~$\mathring I_{\alpha_1}$, so $T^{-1}\circ \I$ is also continuous on~$\mathring I_{\alpha_1}$. Since the discontinuity points of $T^{-1}$ are given by $\{T(\partial I_{\beta}) \mid \beta\in \A {\setminus \{\pi_1^{-1}(1)\}}\}$, so by the maximal length of $I_{\alpha_1}$, non-degenericity of $T$ and the continuity of $T^{-1}\circ \I_I$, we~must have
\begin{equation}\label{symmetric identity}
\I_I(I_{\alpha_1})=T(I_{\alpha_1}).
\end{equation}

By induction on the index $\{\alpha_i, 1\le i\le d\}$, the above identity holds for every $\alpha \in \A$. Because the involution $\I$ reverses the order of $\{I_{\alpha}, \alpha \in \A\}$, that is:
\[
I_{\alpha}<I_{\beta}\implies \I_I(I_{\alpha})>\I_I(I_{\beta}),
\]
{where the order of intervals is given by \eqref{eq: defofintervalorder}, }the identity \eqref{symmetric identity} implies that $T$ also reverses the order of $\{I_{\alpha}, \alpha \in \A\}$, hence $T$ is symmetric.
\end{proof}

\end{document}
